\scanpage{5}
Nécessité claire. Pour la suffisance, on note que les conditions sont stables par passage à une algèbre quotient $B$. Par récurrence noeth, on peut donc supposer que pour tt tel quotient $B=A/J$, avec $J\ne0$, $F_B$ est représentable (par module de t.f.). On peut supposer aussi $A\ne0$. Il existe par \textcircled{\scriptsize iv} un $f\in A$ non nilpotent, tels que $F_{\deleted{\unsure{f}}A_f}$ soit repr. D'autre part, par l'hyp. de récurrence noeth, comme pour tt $n$, on a $f^{n+1}\ne0$, on voit que $F_{A/f^{n+1}A}$ est repr. On peut donc appliquer le lemme 2, qui donne la conclusion.

\noindent\rule{\linewidth}{0.4pt}\par\nobreak\noindent\rule{\linewidth}{0.4pt}

En pratique, la condition (iv) est parfois difficile à vérifier. On va donner d'autres conditions.

\uline{Proposition 4} \quad $A$ \uline{local} \uline{noeth}. Pour que $F$ soit repr. par module de t.f., il f. et s. que $F$ satisfasse les conditions

(i) $F$ exact à gauche

(ii) $F$ commute aux $\varinjlim$

(iii) Pour tt $\widehat A$-module de t.f. $M$,
\[
F(\widehat M)\xrightarrow{\sim}\varprojlim_{n}F(M/{\unsure{\mathfrak m|m}}^{n+1}M),\ \text{et}
\]

(iv) \deleted{\unsure{de}} $F(\widehat M)$ est de t.f. sur $\widehat A$ (il suffit $=$ si $M$ est \uline{fini} sur $A$).

Nécessité claire. Pour suffisance, grâce au Lemme 1 et condition (i), on peut supposer $A$ complet.

% archive sheet number: 4
